% !TEX root = main.tex
% ============================================================================
% PERSIAN ABSTRACT
% ============================================================================

\chapter*{چکیده}
در عصر دیجیتال، تصاویر رنگی به دلیل حجم بالا، ساختار چندکاناله و همبستگی آماری قوی میان پیکسل‌های مجاور، به الگوریتم‌های رمزنگاری تخصصی نیاز دارند که روش‌های سنتی متن‌محور مانند \lr{AES} برای آن‌ها ناکارآمد است. سیستم‌های آشوبی به دلیل حساسیت شدید به شرایط اولیه و تولید دنباله‌های شبه‌تصادفی، گزینه‌ای مناسب برای رمزنگاری تصویر محسوب می‌شوند، اما سیستم‌های کلاسیک مانند نقشه \lr{Logistic} و \lr{Tent} در معرض پدیده «تخریب آشوب» هستند که امنیت الگوریتم را در شرایط خاص کاهش می‌دهد.  در این پژوهش، الگوریتمی پنج‌مرحله‌ای برای رمزنگاری تصاویر رنگی معرفی شده است که از ترکیب سیستم آشوبی سه‌بعدی \lr{Chen} و نه سیستم آشوبی نمایی (\lr{ECM}) بهره می‌برد. در مرحله اول، شرایط اولیه با دقت ۱۵ رقم اعشاری به عنوان کلید رمزنگاری تعریف می‌شوند. در مرحله دوم، پیکسل‌های تصویر با استفاده از دنباله آشوبی حاصل از حل عددی \lr{RK45} سیستم \lr{Chen} با پارامترهای \lr{a}=35، \lr{b}=3 و \lr{c}=28 به‌صورت تصادفی جابه‌جا می‌شوند. در مرحله سوم، برای هر پیکسل به‌صورت مستقل یکی از نه سیستم نمایی شامل \lr{LEL}، \lr{LES}، \lr{LET}، \lr{SEL}، \lr{SES}، \lr{SET}، \lr{TEL}، \lr{TES} و \lr{TET} از طریق یک نقشه \lr{Logistic}  مستقل انتخاب می‌شود. در مرحله چهارم، مقادیر رنگی پیکسل‌ها با عملیات \lr{XOR} با بایت کلید تولیدشده از سیستم نمایی انتخابی ترکیب می‌شوند. در مرحله پنجم، یک لایه \lr{XOR} مستقل اضافی با یک \lr{Logistic} مستقل، عمق رمزنگاری را دوچندان می‌کند.  الگوریتم پیشنهادی بر روی چهار تصویر استاندارد از مجموعه داده \lr{USC-SIPI} در ابعاد ۵۱۲×۵۱۲ و ۲۵۶×۲۵۶ ارزیابی شد. نتایج نشان داد که میانگین آنتروپی شانون تصاویر رمزنگاری‌شده به ۸.۰ بیت (مقدار ایده‌آل) رسیده است. ضریب همبستگی پیکسل‌های مجاور از میانگین ‎۰.۹۹۸۹‎ در تصاویر اصلی به ‎۰.۰۰۸۲‎ در تصاویر رمزنگاری‌شده کاهش یافت که معادل کاهش ‎\lr{۹۹.۲٪}‎ است. نرخ تغییر پیکسل (\lr{NPCR}) برای تصویر \lr{Peppers} به \lr{۹۹.۴۸٪} رسید. معکوس‌پذیری کامل الگوریتم با \lr{MSE} = 0 و \lr{PSNR} = $\infty$ برای تمام تصاویر تأیید شد. فضای کلید الگوریتم از $10^{105}$ تجاوز می‌کند که حمله جستجوی فراگیر را به‌لحاظ محاسباتی ناممکن می‌سازد. نوآوری اصلی این پژوهش، انتخاب پویا و مستقل سیستم آشوبی نمایی برای هر پیکسل است که یک لایه امنیتی منحصربه‌فرد ایجاد می‌کند؛ در صورتی که مهاجم حتی موفق به شناسایی دنباله رمزنگاری یک پیکسل شود، هیچ اطلاعاتی درباره پیکسل‌های دیگر به‌دست نمی‌آورد. مقایسه با روش‌های مرجع نشان می‌دهد الگوریتم پیشنهادی از نظر آنتروپی و \lr{NPCR} رقابت‌پذیر است و با بهینه‌سازی پیاده‌سازی از طریق برداری‌سازی و پردازش موازی \lr{GPU}، پتانسیل دستیابی به سرعت‌های بلادرنگ را داراست.

\textbf{کلمات کلیدی:} رمزنگاری تصویر؛ سیستم آشوبی نمایی؛ سیستم آشوبی \lr{Chen}؛ جایگشت پیکسل؛ آنتروپی شانون؛ \lr{NPCR}؛ \lr{UACI}

\newpage

